% c05 —— 植田一石《ホモロジー的ミラー対称性》第 12–13 节（p019–p021）
% [Tutor Note: 本 chunk 与 c04 的拼接点在 p019 页首；p019 首句承 p018 末句（原文「…を示し／て、」），
%  故本文件从该句中段「…て、M の深谷圏 F(M) が…」处接续；开头补出 c04 句末省略的「示した／证明了」。]

并证明了 $M$ 的 Fukaya 范畴 $\cF(M)$ 可由 $\cF(\pi)$ 描述。这称为 Picard–Lefschetz 理论（Picard–Lefschetz theory）[Sei08]。这个范畴 $\cF(\pi)$ 常被称为 Fukaya–Seidel 范畴，其原型可在 [Kon98] 中看到，而其源头至少可以追溯到 Cecotti–Vafa [CV93] 等弦论学家的工作。Fukaya–Seidel 范畴的定义有若干种流派，其中之一是：以 $\pi$ 的消没闭链的特殊基 $(C_1,\ldots,C_n)$ 为对象集，把它定义为 $\pi$ 的纤维 $M_z:=\pi^{-1}(z)$ 的 Fukaya 范畴 $\cF(M_z)$ 的有向子范畴（directed subcategory）[Sei01b, Sei01a]\footnote{除此之外，还有使用沿纤维分支的全空间的二重覆盖的做法 [Sei08]，以及更近期的 [Sei12a, Seia, Seib] 等。}。在这个定义下，$\cF(\pi)$ 的对象集是有限全序集 $(C_1,\ldots,C_n)$，态射则利用 $\cF(M_z)$ 中的态射由式 (11.1) 给出。
\begin{equation}
  \Hom_{\cF(\pi)}(C_i,C_j) =
  \begin{cases}
    \id_{C_i}, & i=j,\\
    \Hom_{\cF(M_z)}(C_i,C_j), & i<j,\\
    0, & i>j
  \end{cases}
  \tag{11.1}
\end{equation}

由式 (11.1) 容易看出，Fukaya–Seidel 范畴的积 $\mathbf{m}_k$ 在 $k>n$ 时为 $0$，而且即使 $k\le n$ 也只含有有限的信息，因此 Fukaya–Seidel 范畴可以由有限的数据完全描述。

若进一步在 $M_z$ 上赋予 Lefschetz 纤维化的结构，则 $\cF(M_z)$ 的计算便归结为维数再低 2 维的辛流形的 Fukaya 范畴的计算。如此反复，把辛流形的维数降到 2，近全纯圆盘的计数便归结为组合论中「涂色」的问题。

仿射流形（或更一般地，Stein 流形）的 Fukaya 范畴是拓扑几何的（或组合论的）这一想法，可以追溯到余切丛情形的 [Wit95] 以及 [FO97]，并在 [FSS08, FSS09, NZ09, Nad09, AS10a, Abo11, Sei12b] 等中得到发展。Kontsevich [Kon] 通过分析 Seidel 的手法而猜测：任意 Stein 流形 $X$ 都具有一个一般带奇点的 Lagrange 部分流形 $L$，称为 Lagrange 骨架，并且在 $L$ 上存在 dg 范畴的可构造层 $A$，使得 $X$ 的 Fukaya 范畴可作为其整体截面 $H^0(A)$ 得到。与 Kontsevich 这一猜想密切相关的工作有 [DK, STZ14, IU, Nad14, Nad15, NT, Nada, Nadb] 等。

\section{Lefschetz 束}

% [Tutor Note: 原文节标题「Lefschetz 束」在正文中括注为 Lefschetz pencil（代数几何的「束」，即 pencil），
%  与 Lefschetz 纤维化（Lefschetz fibration）并非同一概念。glossary 将「Lefschetz 束」冻结为「Lefschetz 纤维丛」，
%  与此处 pencil 的语境不符，故本节标题按原文保留「束」；原文明写「Lefschetz ファイブレーション」之处，则依
%  glossary 的「〜ファイブレーション → 纤维化」译作「Lefschetz 纤维化」。请审校统一。]

$(M,\omega)$ 为辛流形，$B\subset M$ 为 $M$ 的余维数 2 的辛子流形。称映射 $f\colon M\setminus B\to\PP^1$ 为 Lefschetz 束（Lefschetz pencil），是指：$f$ 仅有有限个非退化临界点，这些临界值互不相同，$B$ 由适当的近全纯局部坐标 $(z_1,\ldots,z_d)$ 定义为 $z_1=z_2=0$，并且在 $B$ 的邻域内 $f$ 由 $(z_1,\ldots,z_d)\mapsto z_1/z_2\in\PP^1$ 给出。因此，沿 $B$ 爆破 $M$，$f$ 便延拓为 $\tilde f\colon\tilde M:=\operatorname{Bl}_M M\to\PP^1$，而 $\tilde f$ 成为 Lefschetz 纤维化。

% [Tutor Note: 原文「\operatorname{Bl}_M M」的足标 M 疑为 B（沿 B 爆破）之笔误，此处按原文保留，请审校确认。]

设 $\pi\colon M\to\C$ 为以 $\C$ 为底空间的 Lefschetz 纤维化，在其临界点集上取定一个全序；$\operatorname{Crit}(f)=(p_1,\ldots,p_n)$。再令临界值集为 $\operatorname{Critv}(f):=f(\operatorname{Crit}(f))=(f(p_1),\ldots,f(p_n))$。在底空间 $\C$ 上取定一个基点 $*$。$\C$ 上满足 $\gamma(0)=*$ 且 $\gamma(1)\in\operatorname{Critv}(f)$ 的单射道路 $\gamma\colon[0,1]\to\C$ 称为消没路径（vanishing path）。对任意消没路径 $\gamma$，其起点处的微分 $\gamma'(0)$ 是 $T_*\C\cong\C$ 的元素。当由消没路径组成的有序集 $(\gamma_1,\ldots,\gamma_n)$ 满足以下条件时，称之为消没路径的临界集（distinguished set of vanishing paths）：

% [Tutor Note: 原文「特集合（distinguished set）」按 glossary 统一作「临界集」；此处的「临界集」意为
%  「特异的、被指定的一组」，与「临界点集」不是同一个意思，特此说明。]

\begin{itemize}
  \item 对任意 $i\in\{1,\ldots,n\}$，有 $\gamma_i(1)=f(p_i)$。
  \item 对任意 $i,j\in\{1,\ldots,n\}$，若 $i\ne j$，则 $\gamma_i([0,1])\cap\gamma_j([0,1])=\{*\}$。
  \item 对任意 $i\in\{1,\ldots,n\}$，有 $\gamma'_i(0)\ne0$，并且若适当选取辐角的分支，则 $\arg\gamma'_1(0)>\cdots>\gamma'_n(0)$。
\end{itemize}

对消没路径 $\gamma\colon[0,1]\to\C$ 与纤维 $M_*:=\pi^{-1}(*)$ 的点 $p\in M_*$，以 $p$ 为起点的 $\gamma$ 的提升 $\tilde\gamma_p\colon[0,1)\to M$ 由下述条件唯一确定：对任意 $t\in[0,1)$，有 $\pi\circ\tilde\gamma_p(t)=\gamma(t)$ 且 $\tilde\gamma'_p(t)\in\left(T_{\tilde\gamma_p(t)}M_{\gamma(t)}\right)^{\perp_\omega}$。这里 $\tilde\gamma'_p(t)\in T_{\tilde\gamma_p(t)}M$ 是 $\tilde\gamma_p$ 在 $t$ 处的微分，而 $\left(T_{\tilde\gamma_p(t)}M_{\gamma(t)}\right)^{\perp_\omega}$ 是 $\gamma(t)\in\C$ 处 $\pi$ 的纤维 $M_{\gamma(t)}:=\pi^{-1}(\gamma(t))$ 在 $\tilde\gamma_p(t)\in M_{\gamma(t)}$ 处的切空间 $T_{\tilde\gamma_p(t)}M_{\gamma(t)}$ 在 $T_{\tilde\gamma_p(t)}M$ 中关于 $\omega$ 的正交补空间。对任意 $t\in[0,1)$，映射 $p\mapsto\tilde\gamma_p(t)$ 给出从 $M_*$ 到 $M_{\gamma(t)}$ 的辛同胚。

利用 $\pi$ 的临界点周围的局部模型 $(z_1,\ldots,z_d)\mapsto z_1^2+\cdots+z_d^2$，可知
\begin{equation}
  V_\gamma := \left\{ p\in M_* \;\middle|\; \lim_{t\to1}\tilde\gamma_p(t)\in\operatorname{Crit}(f) \right\}
  \tag{12.1}
\end{equation}
是与球面同胚的 $M_*$ 的 Lagrange 部分流形。称它为沿 $\gamma$ 的 $\pi$ 的消没闭链（vanishing cycle）。附属于消没路径的临界集 $(\gamma_1,\ldots,\gamma_n)$ 的消没闭链的有序集 $(V_1,\ldots,V_n):=(V_{\gamma_1},\ldots,V_{\gamma_n})$ 称为消没闭链的特殊基（distinguished basis of vanishing paths）。

% [Tutor Note: 原文此处的英文括注作「distinguished basis of vanishing paths」，按上下文应为
%  distinguished basis of vanishing cycles（消没闭链的特殊基）；此处按原文保留，请审校确认。]

对以 $*$ 为起点、绕 $\gamma$ 一周的环路 $\gamma^\#\in\Omega(\C\setminus\operatorname{Critv}(f))$，沿 $\gamma^\#$ 的平行移动给出 $M_*$ 的辛自同胚，称之为沿消没闭链 $V_\gamma$ 的 Dehn 扭转（Dehn twist）。Dehn 扭转在辛映射类群（symplectic mapping class group）$\operatorname{Symp}(M,\omega)/\operatorname{Ham}(M,\omega)$ 中的像只依赖于 $\gamma$ 的同伦类。当 $M_*$ 是 2 维时，上述意义下的 Dehn 扭转在映射类群（mapping class group）$\operatorname{Diff}(M)/\operatorname{Diff}(M)_0$ 中的像就是通常意义的 Dehn 扭转。从辛映射类群到映射类群的映射一般有很大的核 [Sei99, KS02, ST01]。

根据 Donaldson [Don96, Don99] 的著名定理，任意辛流形都具有 Lefschetz 束的结构。此外，在 4 维情形，反过来任意 Lefschetz 束的全空间都容许辛结构，这由 Gompf 所知 [GS99]。给出以亏格 $g$ 的曲面为纤维的 $S^2$ 上的 Lefschetz 纤维化，等价于给出亏格 $g$ 的曲面上的简单闭曲线序列 $(C_1,\ldots,C_n)$，使其在映射类群中满足关系式
\begin{equation}
  \tau_{C_1}\circ\cdots\circ\tau_{C_n}=1
  \tag{12.2}
\end{equation}
关于辛几何中的 Lefschetz 纤维化，有 Donaldson 在 1998 年国际数学家大会上的演讲 [Don98]；关于 Lefschetz 纤维化在 4 维拓扑中的地位，除前述 [GS99] 外，还有 Gompf 的综述 [Gom05] 等。

\section{Fukaya–Seidel 范畴}

对 Lefschetz 纤维化 $\pi\colon M\to\C$，$\C$ 上满足 $\gamma([0,\infty))\cap\operatorname{Critv}(\pi)=\{\gamma(0)\}$、且在紧集之外为平行于负实轴的直线的单射道路 $\gamma\colon[0,\infty)\to\C$，称为以 $-\infty$ 为基点的 $\pi$ 的消没路径。又，把沿消没路径 $\gamma$ 在时刻 $t\in(0,\infty)$ 的消没闭链置为
\begin{equation}
  V_{\gamma,t} := \left\{ p\in M_{\gamma(t)} \;\middle|\; \lim_{s\to0}\tilde\gamma_p(s)\in\operatorname{Crit}(f) \right\}
  \tag{13.1}
\end{equation}
并将这些消没闭链并起来记作
\begin{equation}
  L_\gamma := \bigcup_{t\in[0,\infty)}V_{\gamma,t}
  \tag{13.2}
\end{equation}
称为沿 $\gamma$ 的 Lefschetz 穿透（Lefschetz thimble）。Lefschetz 穿透是与开圆盘同胚的 $M$ 的 Lagrange 部分流形。这里，为方便起见，把由 $\gamma(0)\in\operatorname{Critv}(\pi)$ 上的唯一临界点构成的集合记作 $V_{\gamma,0}$。当由以 $-\infty$ 为基点的消没路径组成的全序集 $(\gamma_i)_{i=1}^n$ 互不相交，而且对充分大的 $t$ 满足 $\operatorname{Im}(\gamma_1(t))>\cdots>\operatorname{Im}(\gamma(n(t))$ 时，称之为以 $-\infty$ 为基点的消没路径的临界集（distinguished set of vanishing paths based at $-\infty$）。此时，对应的 Lefschetz 穿透 $L_i:=L_{\gamma_i}$ 的有序集 $(L_i)_{i=1}^n$ 称为 Lefschetz 穿透的特殊基（distinguished basis of Lefschetz thimbles）。

% [Tutor Note: 原文「Im(γ(n(t))」的括号位置疑为笔误，应为 γ_n(t)；此处按原文保留，请审校确认。]

适当选取并固定包含所有临界值的紧集 $K_1$ 与内部包含 $K_1$ 的紧集 $K_2$。再取充分小的正数 $\varepsilon$，以及 Lefschetz 纤维化的底空间 $\C$ 的微分同胚 $\bar\phi_\varepsilon$，使它在 $K_1$ 上是恒等映射、在 $K_2$ 之外由 $\varepsilon$ 旋转 $z\mapsto\exp(\varepsilon\sqrt{-1})z$ 给出。在 $K_1$ 之外，$\pi$ 成为局部平凡的辛流形族，因此存在 $M$ 的 Hamilton 同胚 $\phi_\varepsilon$，它是 $\bar\phi_\varepsilon$ 的提升。以 $(L_i)_{i=1}^n$ 为对象，
\begin{equation}
  \hom(L_i,L_j) := \mathrm{CF}^{*}(\phi_\varepsilon(L_i),L_j)
  \tag{13.3}
\end{equation}
为态射空间得到的 $A_\infty$ 范畴称为 $\pi$ 的 Fukaya–Seidel 范畴，记作 $\cF(\pi)$。Fukaya–Seidel 范畴中态射的合成，由 $\phi_\varepsilon$ 所诱导的等同
\begin{equation}
  \mathrm{CF}^{*}(\phi_\varepsilon(L_i),L_j)\xrightarrow{\sim}\mathrm{CF}^{*}(\phi_{2\varepsilon}(L_i),\phi_\varepsilon(L_j))
  \tag{13.4}
\end{equation}
与 $E$ 的 Fukaya 范畴中态射的合成
\begin{equation}
  \mathbf{m}_2\colon \mathrm{CF}^{*}(\phi_\varepsilon(L_j),L_k)\otimes\mathrm{CF}^{*}(\phi_{2\varepsilon}(L_i),\phi_\varepsilon(L_j))
  \to\mathrm{CF}^{*}(\phi_{2\varepsilon}(L_i),L_k),
  \tag{13.5}
\end{equation}
以及 Lagrange 交叉 Floer 理论中的连续性映射（continuation map）
\begin{equation}
  \mathrm{CF}^{*}(\phi_{2\varepsilon}(L_i),L_k)\cong\mathrm{CF}^{*}(\phi_\varepsilon(L_i),L_k)
  \tag{13.6}
\end{equation}
复合而成。其中约定 $\phi_{2\varepsilon}:=\phi_\varepsilon\circ\phi_\varepsilon$。$A_\infty$ 范畴中态射的高次合成也类似地定义。容易看出，这个定义与前述把 $\cF(\pi)$ 定义为基点 $*\in\C$ 处 $\pi$ 的纤维 $M_*:=\pi^{-1}(*)$ 的 Fukaya 范畴 $\cF(M)$ 的、附属于消没闭链的特殊基 $(V_i)_{i=1}^N$ 的有向子范畴的定义是一致的。$\cF(\pi)$ 依赖于消没闭链的特殊基以及近复结构等辅助数据，但导出范畴 $D^b\cF(\pi)$ 的拟同构类只依赖于 Lefschetz 纤维化 $\pi$ 的辛同胚类\footnote{Fukaya–Seidel 范畴不依赖于消没闭链特殊基的取法，这是关于 Dehn 扭转的 Lagrange 交叉 Floer 上同调的长正合列 [Sei03] 的推论。}。

% [Tutor Note: 原文「と E の深谷圏における射の合成」中的 E 未加定义；由下段「E ∖ π^{-1}(⋯)」的用法推测，
%  E 指 Lefschetz 纤维化的全空间（即 M）。此处按原文保留 E，请审校确认。]

如果把 Lefschetz 纤维化看作 Morse 函数的类比，那么 Lefschetz 穿透就是下降 Morse 闭链（descending Morse cycle）在 Lefschetz 纤维化中的对应物。取定以 $\C$ 上的点 $*$ 为基点的消没路径的临界集 $(\gamma_i)_{i=1}^n$。$\C\setminus\bigcup_{i=1}^n\gamma([0,1])$ 是可缩的，而且在其上 $\pi$ 成为拓扑平凡的纤维丛，因此 $E\setminus\pi^{-1}(\bigcup_{i=1}^n\gamma([0,1]))$ 与纤维 $M_*:=\pi^{-1}(*)$ 同伦等价。把以 $*$ 为基点的 Lefschetz 穿透粘贴上去，所得空间与 $M$ 同伦等价；
\begin{equation}
  M\sim M_*\cup\bigcup_{i=1}^{n}L_i.
  \tag{13.7}
\end{equation}
因此，Lefschetz 穿透是充当全空间 $M$ 与纤维 $M_*$ 之间的中介的几何对象，Seidel 的纲领就是它在 Fukaya 范畴中的影子。此外，Lefschetz 穿透还有作为超几何级数的积分表示中被积闭链的一面，与量子上同调所伴随的单值保持变形的 Stokes 行列相关的 Dubrovin 猜想 [Dub98] 密切相关。

作为 Lefschetz 纤维化的例子，考虑将 $A_n$ 型奇点的定义多项式 $x^{n+1}$ 用正数 $a$ 扰动而得到的映射
\begin{equation}
  f\colon\C\to\C,\qquad x\mapsto x^{n+1}-ax
  \tag{13.8}
\end{equation}
$f$ 的临界点集为
\begin{align}
  \operatorname{Crit}(f)&:=\{x\in\C\mid (n+1)x^n=a\}
  =\left\{\zeta_n^ix_0\right\}_{i=0}^{n-1},\notag\\
  &\hspace{3em}x_0:=\sqrt[n]{a/(n+1)},\qquad
  \zeta_n:=\exp(2\pi\sqrt{-1}/n)
  \tag{13.9}
\end{align}
而临界值
\begin{equation}
  x^{n+1}-ax=\frac{a}{n+1}x-ax=-\frac{n}{n+1}ax
  \tag{13.10}
\end{equation}
分布在圆周上。把消没路径的临界集取成如图 13.1 那样，则对应的 Lefschetz 穿透就如图 13.2 那样。它们满足
\begin{equation}
  \#\left(\phi_\varepsilon(L_i)\cap L_j\right)=
  \begin{cases}
    1, & j=i\ \text{或}\ j=i+1,\\
    0, & \text{其他}
  \end{cases}
  \tag{13.11}
\end{equation}
因此，适当选取分次后，Fukaya–Seidel 范畴中的态射空间由
\begin{equation}
  \hom^k(L_i,L_j)=
  \begin{cases}
    \C\cdot\id_{L_i}, & j=i\ \text{且}\ k=0,\\
    \C, & j=i+1\ \text{且}\ k=1,\\
    0, & \text{其他}
  \end{cases}
  \tag{13.12}
\end{equation}
给出。

% CHUNK c05 DONE pages p019–p021
